\chapter{Blood Glucose}
\section*{What Is This Test?}
The blood glucose test measures the concentration of glucose in the blood. Glucose is the primary energy substrate for cellular metabolism, particularly for the central nervous system (which is obligatorily dependent on glucose except during prolonged fasting/ketosis). Blood glucose is tightly regulated by the opposing hormones insulin (promoting glucose uptake, glycogen synthesis, and glycolysis) and glucagon, cortisol, epinephrine, and growth hormone (promoting glycogenolysis, gluconeogenesis, and lipolysis). Blood glucose is measured as part of fasting plasma glucose (FPG), random glucose, 2-hour postprandial glucose, or as the 2-hour value in a 75-gram oral glucose tolerance test (OGTT). The diagnosis of diabetes mellitus (DM) is established by any of the following: (1) FPG $\geq$126 mg/dL on two occasions, (2) 2-hour OGTT plasma glucose $\geq$200 mg/dL, (3) HbA1c $\geq$6.5\% (per NGSP-certified and DCCT-standardized assay), or (4) random plasma glucose $\geq$200 mg/dL with symptoms of hyperglycemia (polyuria, polydipsia, unexplained weight loss). Prediabetes is defined by: FPG 100--125 mg/dL (impaired fasting glucose, IFG), 2-hour OGTT glucose 140--199 mg/dL (impaired glucose tolerance, IGT), or HbA1c 5.7--6.4\%. Blood glucose is measured by the glucose oxidase method or hexokinase method.
\section*{Alternative Names}
Blood Sugar; Fasting Blood Glucose (FBG); Fasting Plasma Glucose (FPG); Random Blood Glucose; Serum Glucose; Plasma Glucose
\section*{Principle}
Glucose is measured by enzymatic colorimetric methods. Glucose oxidase method: Glucose + O2 $\rightarrow$ gluconic acid + H2O2 (glucose oxidase); H2O2 + chromogenic oxygen acceptor (4-aminoantipyrine) $\rightarrow$ colored product (peroxidase). The hexokinase method (reference method): Glucose + ATP $\rightarrow$ glucose-6-phosphate + ADP (hexokinase); glucose-6-phosphate + NADP+ $\rightarrow$ 6-phosphogluconate + NADPH (glucose-6-phosphate dehydrogenase); NADPH measured spectrophotometrically at 340 nm. The hexokinase method is more specific than glucose oxidase. Specimen: plasma from a sodium fluoride/potassium oxalate (gray-top) tube (fluoride inhibits glycolysis in vitro; without fluoride, glucose decreases 5--7\% per hour at room temperature due to continued glycolysis by RBCs and WBCs). Capillary blood glucose (fingerstick glucometer) uses glucose oxidase or glucose dehydrogenase test strips.
\section*{History}
In 1921, Frederick Banting and Charles Best isolated insulin from pancreatic extracts, and its first successful clinical use in 1922 transformed diabetes from a rapidly fatal disease into a manageable one; Banting and Macleod received the Nobel Prize in 1923. The development of glucose oxidase test strips (Dextrostix, 1964) and of home blood glucose meters in the 1970s--1980s allowed patients to self-monitor glucose. Diagnostic thresholds were standardized by the National Diabetes Data Group (1979) and the World Health Organization (1980), and in 1997 the American Diabetes Association lowered the fasting diagnostic threshold to $\geq$126 mg/dL. The discovery of glycated hemoglobin (HbA1c) in the late 1960s, validated for outcome prediction by the Diabetes Control and Complications Trial (1993), provided a long-term index of glycemic control, and HbA1c was formally added to the diagnostic criteria in 2010.
\section*{Why This Test Is Ordered}
\begin{itemize}
  \item Diagnosis of diabetes mellitus and prediabetes using fasting plasma glucose, random glucose, or the 2-hour OGTT
  \item Screening of asymptomatic adults who are overweight (BMI $\geq$25) or have other risk factors
  \item Evaluation of symptomatic hyperglycemia (polyuria, polydipsia, unexplained weight loss)
  \item Investigation of suspected hypoglycemia, confirming the Whipple triad
  \item Routine screening for gestational diabetes in pregnancy
  \item Monitoring of glycemic control in patients with established diabetes, including at the point of care
  \item Preoperative assessment and monitoring of the critically ill, where stress hyperglycemia is common
\end{itemize}
\section*{Is This a Primary or Secondary Test}
Blood glucose is the primary diagnostic test for diabetes mellitus and hypoglycemia, and fasting plasma glucose is the most commonly used first-line screening test in at-risk adults. It is also the central monitoring test in established diabetes, where it is complemented by HbA1c. In pregnancy, glucose measurement is used routinely to screen for gestational diabetes, typically as part of an oral glucose tolerance test.
\section*{How This Test Is Performed}
A blood sample is collected by standard venipuncture into a sodium fluoride/potassium oxalate (gray-top) tube when a fasting glucose is ordered, since fluoride inhibits in vitro glycolysis; a serum separator tube is acceptable if the sample is centrifuged and processed promptly. Glucose is measured by the glucose oxidase or hexokinase enzymatic method on an automated chemistry analyzer. Capillary blood glucose measured on a bedside glucometer is used for monitoring but is not recommended for the formal diagnosis of diabetes. Results are typically available within a few hours, or within minutes at the point of care.
\section*{What Preparation Is Needed}
Fasting plasma glucose requires no caloric intake for 8 hours, usually overnight, with water permitted. The 2-hour oral glucose tolerance test requires an 8--12 hour fast and ingestion of a 75-gram glucose load. Random glucose requires no preparation. Patients should inform their clinician of medications that raise glucose (e.g., corticosteroids, thiazide diuretics, beta-blockers), but medications are generally not held unless instructed.
\section*{Contraindications and Precautions}
There are no absolute contraindications to venipuncture. Precautions: acute stress hyperglycemia in hospitalized or critically ill patients is common and should not be over-interpreted as diabetes without confirmatory testing after recovery; a sodium fluoride tube should be used when sample processing is delayed, since glucose decreases 5--7\% per hour at room temperature in plain tubes; the diagnosis of diabetes should be confirmed by a second test on a separate occasion, except in the setting of unequivocal symptomatic hyperglycemia; and critical values (glucose <40--50 mg/dL or >400--500 mg/dL) require immediate clinical notification and treatment.
\section*{Risks}
Minimal risk. Slight pain or bruising at the venipuncture site. Rarely: hematoma, infection, or vasovagal syncope.
\section*{What Happens After the Test}
Results are typically available within a few hours and are interpreted against the diagnostic thresholds for diabetes and prediabetes. An abnormal fasting or OGTT value should be confirmed on a second occasion before a diagnosis of diabetes is established. Patients with new diabetes are referred for endocrine evaluation, diabetes self-management education, and cardiovascular risk reduction; symptomatic hypoglycemia is treated immediately with oral glucose or intravenous dextrose.
\section*{Understanding Results}
\subsection*{Hypoglycemia (<70 mg/dL)}
Whipple triad: symptoms of hypoglycemia + low plasma glucose + resolution of symptoms with glucose administration. Causes: exogenous insulin or sulfonylurea excess, insulinoma, alcohol binge with fasting, liver failure, adrenal insufficiency, hypothyroidism, sepsis. Neuroglycopenic symptoms: confusion, altered behavior, seizures, coma. Autonomic symptoms: diaphoresis, palpitations, tremor, anxiety, hunger.
\subsection*{Hyperglycemia}
FPG 100--125 mg/dL: impaired fasting glucose (prediabetes). FPG $\geq$126 mg/dL: diabetes mellitus (on two occasions). Random glucose $\geq$200 mg/dL with symptoms: diabetes. 2-hour OGTT $\geq$200 mg/dL: diabetes. Stress hyperglycemia (critical illness, sepsis, myocardial infarction, surgery). Diabetic ketoacidosis (DKA): hyperglycemia >250 mg/dL, anion gap metabolic acidosis, ketonemia/ketonuria. Hyperosmolar hyperglycemic state (HHS): glucose >600 mg/dL, hyperosmolarity >320 mOsm/kg, no significant ketosis or acidosis.
\section*{Normal Results}
Fasting plasma glucose: 70--99 mg/dL (3.9--5.5 mmol/L). 2-hour postprandial or OGTT glucose: <140 mg/dL (<7.8 mmol/L). Random glucose: <140 mg/dL (normally <200 mg/dL). HbA1c: <5.7\%. Critical values: hypoglycemia <40--50 mg/dL (immediate treatment required); severe hyperglycemia >400--500 mg/dL (risk of DKA or HHS).
\section*{Follow-up Tests}
HbA1c, oral glucose tolerance test (OGTT), serum insulin and C-peptide (to differentiate type 1 from type 2), anti-GAD65, anti-IA2, anti-insulin, and anti-ZnT8 antibodies (for autoimmune type 1 diabetes), basic metabolic panel (electrolytes, anion gap for DKA), serum and urine ketones, and serum osmolality (for HHS).
\section*{References}
\begin{enumerate}
  \item American Diabetes Association. Classification and diagnosis of diabetes: Standards of Medical Care in Diabetes --- 2024. Diabetes Care. 2024;47(Suppl 1):S20--S42.
  \item Sacks DB, Arnold M, Bakris GL, et al. Guidelines and recommendations for laboratory analysis in the diagnosis and management of diabetes mellitus. Clin Chem. 2011;57(6):793--798.
  \item Banting FG, Best CH. The internal secretion of the pancreas. J Lab Clin Med. 1922;7(5):251--266.
  \item World Health Organization. Definition, Diagnosis and Classification of Diabetes Mellitus and its Complications: Report of a WHO Consultation. Geneva: World Health Organization; 1999.
  \item Sacks DB. A1C versus glucose testing: a comparison. Diabetes Care. 2011;34(2):518--523.
\end{enumerate}
\clearpage
